%=============================================================================!
% 6.0  CHAPTER OPENING                                                        !
%=============================================================================!
Newton's laws require the total force on each body, but a constraint
force is often known only through the motion it must maintain. The push
of a wire keeps a bead on its curve, and the force of a pendulum's rod
keeps the bob at a fixed distance from the pivot. Their magnitudes depend
on the motion we are trying to find. When these forces do no work, we can
instead describe the allowed motion with independent coordinates and
write the kinetic and potential energies in those coordinates.

This leads to the Lagrangian formulation of mechanics. Hamilton's
principle compares whole paths through time, and the Euler--Lagrange
equations turn that comparison into equations of motion. We apply them
to the bead on the wire of Chapter~3, then use the dependence of the
Lagrangian on its coordinates to identify conserved momenta. When the
constraint force itself is needed, a Lagrange multiplier supplies it;
Chapter~11 uses this method to hold molecular bonds at fixed lengths.
The conjugate momenta introduced here also prepare the Hamiltonian
formulation in Chapter~7.
